\section{Evaluation methodology}\label{sec:eval}

The empirical scope covers implementation correctness, network resource costs, evidence preservation and a small retrospective metadata-classification diagnostic. Fixed-arrival mechanism simulations are reported separately from network measurements. The energy model and budgeted attacker-selection rule are theoretical design components. We do not test a privacy-per-joule superiority hypothesis, a server--edge comparison or a 10\,Gbps throughput claim.

\subsection{Testbed and workloads}\label{sec:testbed}
We project the public HIGGS dataset~\cite{whiteson2014higgs} to its binary label and first two input features. The source contains 11 million simulated collision records and is distributed under CC BY 4.0. The projection, site-partition order, question banks and references have separate manifests and content hashes. This reduced analytics workload is not a HIGGS classification benchmark. Four query families are used: aggregate statistics, group-by, record count and a small weighted federated-training workflow. For the latter the source bank specifies an initial step; our declared deterministic controller performs three updates, feeding the aggregate weights into the next step. We do not attribute this extension to autonomous LLM replanning.

The banks have question-, template- and cohort-separated TRAIN, SELECT and reserved final-test partitions. The calibration uses one question from each of four family and three requested-size strata, with two repetitions. Repetitions of the same question are not independent new questions. The final-test partition is not executed in this engineering study. Independent direct computation checks every aggregate update and final tool result against the same query definition without the network path.

The Mac Studio experiment runs 20 logical partitions and persistent TLS connections on one machine. Its minimal TLS RPC baseline uses the same query handler without the all-site scheduling controls. The shaped unshaped-mode baseline retains those controls. The four-instance deployment uses four \nolinkurl{m6i.2xlarge} EC2 hosts in one availability zone: T1 is the client, T2 and T3 are relays or NetShaper middleboxes, and T4 holds all 20 logical partitions. The data path uses private IP addresses. Captures at T2 and T3 observe the same TCP or QUIC/UDP segment; SSH management and evidence replication are outside that segment. This is one multiplexed observer link and one physical data endpoint, not 20 independently observed physical sites.

\subsection{Configuration selection and comparison}
The Mac TRAIN grid comprises minimal TLS, unshaped with all-site controls, fixed-1024-byte padding, bucketed padding, rounds, and DP/combined configurations with bias multipliers in $\{0,1,2,4,8\}$. The DP arrival variant uses $H=10$\,s, $\tau=0.5$\,s, a per-direction input cap of 8192 bytes and $\Delta=16384$ bytes. For $N=20$, $\varepsilon_{\mathrm{dir}}=0.2$ gives a declared session accounting budget of 8 over the 40 directions. A bias multiplier $m$ denotes $m\Delta/\varepsilon_{\mathrm{dir}}$. Rounds use $W=0.5$\,s, cover probability 1 and a 1024-byte reply frame.

We freeze the five baselines and the two cheapest fully passing TRAIN candidates from each DP family, then evaluate that shortlist on SELECT. The selected candidate in each family minimises mean framed plaintext bytes among candidates passing every SELECT trial. Truncated or incomplete trials remain failures, and host scheduling failures halt acquisition. This rule is functional/cost calibration, not \cref{alg:controller}: it uses neither an attacker score nor a calibrated joule budget. The cloud run freezes DP bias 1 and combined bias 2 from this process and evaluates seven modes: unshaped with controls, padding, rounds, DP, combined, and NetShaper with shaping disabled or enabled. It executes 12 TRAIN questions twice per mode, in a fixed seeded permutation of mode phases. We do not pool this three-step campaign with the older one-step pilot.

NetShaper~\cite{sabzi2024netshaper} runs in two middleboxes with the same client and server tools. Its resource configuration is preserved with the artifact. We do not match its privacy accounting to the tool-layer DP mechanism, so differences in byte cost cannot establish privacy superiority. The separate minimal TLS baseline is measured on Mac, not represented as an additional cloud mode.

\subsection{Model compatibility checks}
Two existing local models, \nolinkurl{mistral-small3.1:24b} and \nolinkurl{qwen2.5:14b-instruct-q4_K_M}, generate typed initial plans and structured numerical summaries. Decoding uses temperature 0, seed 20261021, an 8192-token context and at most 2048 predicted tokens. The accepted plan must match the question's reference; no fallback plan is substituted. After tool execution, the model receives the verified aggregate and must return its numerical JSON contents. Integers are exact; finite real values use relative tolerance $10^{-6}$ and absolute tolerance $10^{-8}$. Each model runs 12 TRAIN and 12 SELECT questions, with source, model and input contracts frozen before SELECT. Prompts and raw responses are retained. This measures structured-result consistency, not open-ended factual quality. The cloud workload uses a deterministic scripted controller; there is no hosted-model or LangGraph end-to-end evaluation.

\subsection{Public-contract TLS follow-up}\label{sec:localcontractmethod}
A separate exploratory Mac experiment uses three loopback TLS endpoints and generated integer tables, not the HIGGS bank or an LLM. Twelve TRAIN jobs fix planning and compute bounds at 0.04 and 0.08\,s; each bound rounds twice the observed maximum plus 0.005\,s up to a 0.04\,s multiple. The common contract has three rounds, 512/1024-byte call/reply frames and a 0.6\,s horizon. Fresh SELECT jobs use 1024 or 4096 rows per site, versus 512 on TRAIN, with scripted feedback-dependent routing. Family and round count are coupled in this pilot. Within-job policy order is randomised. The final frequency comparison freezes six or ten periodic ticks through the same horizon before dispatch. Planning/compute delays include controlled perturbations; they do not calibrate LLM bounds. Full-lifecycle loopback captures include setup, controls, ACKs and teardown. Paired-job bootstrap intervals are descriptive, conditional on this workload and host.

\subsection{Observer and privacy scope}\label{sec:attack}
The theoretical observer in \cref{sec:model} sees separate site links. The EC2 capture instead sees a multiplexed tunnel. We do not assign outer packets to logical site identities or claim measured participation privacy from that trace. A future attacker study would need observations matching the claimed threat model, retraining under each configuration, question-disjoint training and selection, a frozen final test, and question-level uncertainty estimates. The retrospective classifier in \cref{sec:retrospective} is limited to workload family and requested size on the multiplexed link; it is not a confirmatory final-test study, per-site participation attack, ROC analysis or empirical DP certification. Passing the transport and reference checks does not establish low leakage; the formal results remain conditional on their stated mechanism and deployment assumptions.

\subsection{Network metrics and evidence integrity}
Packet sizes are IPv4 lengths including TCP/UDP headers and ACKs, excluding Ethernet and the other deployment links. We use half-open per-session windows fixed by client timestamps, and separately report the full phase and bytes outside the session windows. Tool latency is measured from session start to the final tool result, including the configured 1\,s public start lead; session completion also includes any fixed covert tail. The two metrics must not be conflated with LLM answer latency or steady-state throughput.

Each capture has explicit ready and stopped timestamps and kernel-drop counters. NIC offloads are disabled for the matched capture and restored afterwards; host resources and clock offsets are screened before phases. Packet fingerprints compare both observations as an integrity diagnostic; a boundary ACK discrepancy is reported rather than silently removed. Each session has an fsync-backed intent/commit ledger and synchronous evidence replication to T4; phase results and captures are copied to Mac. Checksums, a separate controller copy and an independent restore drill protect against loss or accidental modification. Failed and superseded runs are retained under separate run identifiers, without automatic replay.

\subsection{Energy scope}\label{sec:energy}
No instance exposes an energy counter in the completed cloud testbed, and no external power meter or edge device is used. Consequently we report no measured joules, fitted per-node energy coefficients, privacy gain per joule or CO$_2$ estimate. Bytes and latency describe resource demand and do not substitute for physical energy. Calibrating \cref{eq:energy} with suitable package/DRAM counters or external instrumentation, including idle subtraction and device-specific uncertainty, remains future work~\cite{linux2026powercap}.

\subsection{Reproducibility}
The artifact records source and runtime hashes, the HIGGS projection and partition manifests, question and reference hashes, selected parameters, prompts, model responses, raw captures, per-session ledgers and cleanup records. Private SSH/TLS keys are excluded from the evidence release. The scientific implementation and public dataset licence are distinct: a release licence and public archive identifier must be supplied by the authors before submission. The completed study is descriptive; two repetitions and 12 selected questions per cloud mode do not support broad population-level or deployment-independent performance claims.
